\chapter{画面の見方}\label{chap:screen}

\section{起動したところ}

\LS を起動すると、次のような画面が出ます。ウインドウは大きく\textbf{左・中央・右}の3つに分かれていて、下に工程のボタンが並びます。

\screen{win_empty}{起動した直後の画面}{%
  \mk{queue}{1}\mk[e]{addbutton}{2}\mk{graph}{3}\mkd[isw]{preview}{4}
  \mk[w]{presets}{5}\mk[w]{tabs}{6}\mkd[w]{inspector}{7}\mk[nw]{process}{8}\mk[w]{status}{9}
  \drag{(0.47,0.60)}{($(queue-e)+(0.004,0)$)}
  \node[lscallout] at (0.47,0.66) {Finder から動画やフォルダを\\ここへドラッグ＆ドロップ};}

\begin{marklist}
  \item \textbf{入力キュー}：処理するファイルの一覧です。ここへファイルやフォルダを\textbf{ドラッグ＆ドロップ}して追加できます。
  \item \ui{追加…}：ファイルを選んで追加します（メニューの「ファイル → 開く…」\key{⌘}\key{O} と同じ）。
        \ui{削除} は選んだファイルを一覧から外し、\ui{クリア} は一覧・結果・設定をすべて初期状態に戻します（元のファイルは消えません）。
  \item \textbf{品質グラフ}：品質評価をすると、フレームごとの点数がここにグラフで出ます。
  \item \textbf{プレビュー}：いま扱っている画像を大きく表示します。
  \item \textbf{プリセット}：すべての設定をひとまとめにして保存・読み込みします（第\ref{sec:preset}節）。
  \item \textbf{工程タブ}：右側の設定を「品質評価・アライメント・スタック・仕上げ・出力」の4つに分けています。
        工程が終わると次のタブへ自動で切り替わります。
  \item \textbf{設定パネル}：いまのタブの設定が並びます。上下にスクロールできます。
  \item \textbf{工程のボタン}：\ui{品質評価}\ui{アライメント}\ui{スタック}\ui{書き出し…}。前の工程が終わるまでは押せません（灰色）。
  \item \textbf{状態表示}：いま何をしているか、次に何をすればよいかが出ます。処理中は進み具合と残り時間、\ui{中断} ボタンが出ます。
\end{marklist}

\section{処理したあとの画面}

スタックと仕上げまで終えると、画面の各部にいろいろな情報が出ます。

\screen{win_overview}{スタックと仕上げを終えた画面}{%
  \mk{queue}{1}\mk{graph}{2}\mkd[inw]{preview}{3}\mk[nw]{viewmode}{4}\mk[s]{frameslider}{5}\mk[sw]{zoom}{6}\mk[s]{display}{7}
  \mk[nw]{banner}{8}\mk[w]{tabs}{9}\mkd[w]{inspector}{10}\mk[s]{process}{11}\mk[w]{status}{12}}

\begin{marklist}
  \item 入力キューの各行には、ファイル名と「フレーム数・画像の大きさ・形式」、処理の段階（\uil{未処理}\uil{品質評価済}\uil{アライメント済}\uil{完了}）が出ます。
  \item 品質グラフ：オレンジの横線は「参照画像に使う上位の割合」です（第\ref{sec:qs-quality}節）。
  \item プレビュー：スタックの結果（仕上げを掛けたもの）が出ています。
  \item 表示の切り替え：\uil{フレーム}（撮った1コマ）、\uil{参照}（位置合わせのお手本）、\uil{結果}（スタックの結果）。
  \item フレーム送り：スライダーと \ui{←}\ui{→} で、表示するフレームを選びます。
  \item 拡大率：\uil{全体}（ウインドウに合わせる）、\uil{100\%}（画面の画素と1対1）、\uil{200\%}、\uil{400\%}。
  \item \ui{表示} メニューと \uil{配置を編集}：位置合わせ領域の枠の表示、品質による色分け、画面だけ明るくする表示を切り替えます。
  \item 警告の帯：注意すべきことがあると、プレビューの上に出ます。\ui{×} で閉じられます（第\ref{chap:trouble}章）。
  \item 工程タブ：このときは「仕上げ・出力」が選ばれています。
  \item 設定パネル：仕上げの設定が並んでいます。
  \item 工程のボタン：一度実行すると \ui{品質を再評価}\ui{再アライメント}\ui{再スタック} に名前が変わります。
  \item 状態表示：最後の処理の結果（ここでは RGB のずれの量）が出ます。
\end{marklist}

\section{プレビューの下の道具}\label{sec:toolbar}

\screen{toolbar}{プレビューの下の道具}{%
  \mk[n]{viewmode}{1}\mk[n]{order}{2}\mk[n]{slider}{3}\mk[n]{step}{4}\mk[n]{info}{5}
  \mk[s]{zoom}{6}\mk[s]{display}{7}\mk[s]{apedit}{8}\mk[s]{apcount}{9}}

\begin{marklist}
  \item \textbf{表示の切り替え}：\uil{フレーム}・\uil{参照}・\uil{結果}。参照は品質評価のあと、結果はスタックのあとに選べます。
  \item \textbf{並び順}：スライダーでフレームを送る順番。\uil{時系列}（撮った順）と\uil{品質順}（写りのよい順）があります。
        品質評価のあとは品質順にすると、左端が最良、右へ行くほど悪いフレームになり、見比べやすくなります。
  \item \textbf{フレームのスライダー}：ドラッグしてフレームを選びます。
  \item \textbf{コマ送り}：\ui{←}\ui{→} で1枚ずつ送ります。押し続けると連続して送ります。キーボードの \key{←}\key{→} も使え、
        \key{shift} を押しながらだと10枚、\key{option} を押しながらだと100枚ずつ進みます。
  \item \textbf{表示中の画像の情報}：フレームなら「\#番号 / 総数・上位何\%・品質」、ほかに「参照画像」「スタック結果」など。
  \item \textbf{拡大率}：\key{⌘} を押しながらマウスのホイール、またはトラックパッドのピンチでも拡大・縮小できます（カーソルの位置が中心）。
        拡大しているときは2本指のスクロール（ホイール）で見る場所を動かせます。
  \item \ui{表示} メニュー：\uil{位置合わせ領域を表示}（枠の表示）、\uil{品質で色分け}（第\ref{sec:heatmap}節）、
        \uil{表示を明るくする（保存には影響しません）}（暗い画像を画面の上でだけ明るく見せる）。
  \item \uil{配置を編集}：オンにすると、プレビューをクリックして位置合わせ領域を足したり、選んで \key{delete} で消したりできます（第\ref{sec:apedit}節）。
  \item 位置合わせ領域の数と大きさ（例：\uil{位置合わせ領域 110 個 / 48 px}）。
\end{marklist}

\begin{tip}[画素の値を見る]
  プレビューの上にマウスを置くと、プレビューの下の道具の右端（位置合わせ領域の数の右）に、カーソル位置の座標と画素の値（0〜1）が出ます。
  白飛び（1.0 に張り付く）や黒つぶれを確かめるのに便利です。
\end{tip}

\section{設定パネルの使い方}

右の設定パネルは、\textbf{見出し}（\uil{▼ 各フレームの品質評価} のような太字の行）ごとにまとまっています。
見出しをクリックすると、その欄を\textbf{たたむ}（▶）・\textbf{開く}（▼）を切り替えられます。
ふだん触らない設定（\uil{詳細設定}など）はたたんでおくと見やすくなります。開閉の状態は次に起動したときも保たれます。

\begin{caution}[起動するたびに設定は初期値に戻ります]
  \LS は、起動するたびに入力キュー・設定・仕上げをすべて初期値から始めます（同じ条件を毎回確実に再現するため）。
  よく使う設定は\textbf{プリセット}に保存しておき、読み込んで使ってください（第\ref{sec:preset}節）。
  なお、解析済みのファイルを開いたときは、解析結果を使い回すのに必要な設定だけはサイドカーから戻ります。
\end{caution}

\section{左右の欄を隠す}

プレビューを広く使いたいときは、メニューの「表示」から左右の欄を隠せます。\key{⌘}\key{1} で左（入力キューと品質グラフ）、
\key{⌘}\key{2} で右（設定パネル）を、隠す・出すを切り替えます。
